% =============================================================================
% paper-to-youtube 한국어 프리앰블 (sustech 테마 재사용, 16:9, 학교 로고·크레딧 제거)
% main.tex 맨 위에서  % =============================================================================
% paper-to-youtube 한국어 프리앰블 (sustech 테마 재사용, 16:9, 학교 로고·크레딧 제거)
% main.tex 맨 위에서  % =============================================================================
% paper-to-youtube 한국어 프리앰블 (sustech 테마 재사용, 16:9, 학교 로고·크레딧 제거)
% main.tex 맨 위에서  % =============================================================================
% paper-to-youtube 한국어 프리앰블 (sustech 테마 재사용, 16:9, 학교 로고·크레딧 제거)
% main.tex 맨 위에서  \input{preamble_ko}  한 줄로 불러온다.
% 빌드: latexmk -xelatex main.tex   (latexmkrc 가 ./sustech-theme// 를 TEXINPUTS 에 추가)
% =============================================================================
\documentclass[aspectratio=169,10pt]{beamer}
\usepackage{fontspec}
% 한글 폰트: Windows 맑은 고딕 → 없으면 Noto Sans CJK KR → 나눔고딕
\IfFontExistsTF{Malgun Gothic}{%
  \setsansfont{Malgun Gothic}[AutoFakeSlant=0.15]\setmainfont{Malgun Gothic}[AutoFakeSlant=0.15]%
}{\IfFontExistsTF{Noto Sans CJK KR}{%
  \setsansfont{Noto Sans CJK KR}[AutoFakeSlant=0.15]\setmainfont{Noto Sans CJK KR}[AutoFakeSlant=0.15]%
}{%
  \setsansfont{NanumGothic}[AutoFakeSlant=0.15]\setmainfont{NanumGothic}[AutoFakeSlant=0.15]%
}}
\usetheme{sustech}

\usepackage{amssymb}
\usepackage{tikz}
\usetikzlibrary{positioning,arrows.meta,shapes}
\graphicspath{{figures/}}

% 조건 플래그(테마의 시각화 구분 페이지용, 사용 안 함)
\newif\ifvizbgdiv
\vizbgdivfalse

% --- 중국어 고정 문구를 한국어로 -------------------------------------------
\renewcommand{\figureprefix}{그림}
\renewenvironment{callout}[1][핵심]{\begin{alertblock}{#1}}{\end{alertblock}}
\makeatletter
\renewcommand{\titlemetaline}{%
  \ifdefempty{\presenter}{}{\presenter\metasep}%
  \ifdefempty{\presvenue}{}{\presvenue\metasep}%
  \insertdate}
\makeatother

% --- 브랜딩: 학교 로고·크레딧 끔 (채널 로고가 있으면 \setlogo[0.14\paperheight]{경로}) ---
\setlogo{}
\hidecredits
  한 줄로 불러온다.
% 빌드: latexmk -xelatex main.tex   (latexmkrc 가 ./sustech-theme// 를 TEXINPUTS 에 추가)
% =============================================================================
\documentclass[aspectratio=169,10pt]{beamer}
\usepackage{fontspec}
% 한글 폰트: Windows 맑은 고딕 → 없으면 Noto Sans CJK KR → 나눔고딕
\IfFontExistsTF{Malgun Gothic}{%
  \setsansfont{Malgun Gothic}[AutoFakeSlant=0.15]\setmainfont{Malgun Gothic}[AutoFakeSlant=0.15]%
}{\IfFontExistsTF{Noto Sans CJK KR}{%
  \setsansfont{Noto Sans CJK KR}[AutoFakeSlant=0.15]\setmainfont{Noto Sans CJK KR}[AutoFakeSlant=0.15]%
}{%
  \setsansfont{NanumGothic}[AutoFakeSlant=0.15]\setmainfont{NanumGothic}[AutoFakeSlant=0.15]%
}}
\usetheme{sustech}

\usepackage{amssymb}
\usepackage{tikz}
\usetikzlibrary{positioning,arrows.meta,shapes}
\graphicspath{{figures/}}

% 조건 플래그(테마의 시각화 구분 페이지용, 사용 안 함)
\newif\ifvizbgdiv
\vizbgdivfalse

% --- 중국어 고정 문구를 한국어로 -------------------------------------------
\renewcommand{\figureprefix}{그림}
\renewenvironment{callout}[1][핵심]{\begin{alertblock}{#1}}{\end{alertblock}}
\makeatletter
\renewcommand{\titlemetaline}{%
  \ifdefempty{\presenter}{}{\presenter\metasep}%
  \ifdefempty{\presvenue}{}{\presvenue\metasep}%
  \insertdate}
\makeatother

% --- 브랜딩: 학교 로고·크레딧 끔 (채널 로고가 있으면 \setlogo[0.14\paperheight]{경로}) ---
\setlogo{}
\hidecredits
  한 줄로 불러온다.
% 빌드: latexmk -xelatex main.tex   (latexmkrc 가 ./sustech-theme// 를 TEXINPUTS 에 추가)
% =============================================================================
\documentclass[aspectratio=169,10pt]{beamer}
\usepackage{fontspec}
% 한글 폰트: Windows 맑은 고딕 → 없으면 Noto Sans CJK KR → 나눔고딕
\IfFontExistsTF{Malgun Gothic}{%
  \setsansfont{Malgun Gothic}[AutoFakeSlant=0.15]\setmainfont{Malgun Gothic}[AutoFakeSlant=0.15]%
}{\IfFontExistsTF{Noto Sans CJK KR}{%
  \setsansfont{Noto Sans CJK KR}[AutoFakeSlant=0.15]\setmainfont{Noto Sans CJK KR}[AutoFakeSlant=0.15]%
}{%
  \setsansfont{NanumGothic}[AutoFakeSlant=0.15]\setmainfont{NanumGothic}[AutoFakeSlant=0.15]%
}}
\usetheme{sustech}

\usepackage{amssymb}
\usepackage{tikz}
\usetikzlibrary{positioning,arrows.meta,shapes}
\graphicspath{{figures/}}

% 조건 플래그(테마의 시각화 구분 페이지용, 사용 안 함)
\newif\ifvizbgdiv
\vizbgdivfalse

% --- 중국어 고정 문구를 한국어로 -------------------------------------------
\renewcommand{\figureprefix}{그림}
\renewenvironment{callout}[1][핵심]{\begin{alertblock}{#1}}{\end{alertblock}}
\makeatletter
\renewcommand{\titlemetaline}{%
  \ifdefempty{\presenter}{}{\presenter\metasep}%
  \ifdefempty{\presvenue}{}{\presvenue\metasep}%
  \insertdate}
\makeatother

% --- 브랜딩: 학교 로고·크레딧 끔 (채널 로고가 있으면 \setlogo[0.14\paperheight]{경로}) ---
\setlogo{}
\hidecredits
  한 줄로 불러온다.
% 빌드: latexmk -xelatex main.tex   (latexmkrc 가 ./sustech-theme// 를 TEXINPUTS 에 추가)
% =============================================================================
\documentclass[aspectratio=169,10pt]{beamer}
\usepackage{fontspec}
% 한글 폰트: Windows 맑은 고딕 → 없으면 Noto Sans CJK KR → 나눔고딕
\IfFontExistsTF{Malgun Gothic}{%
  \setsansfont{Malgun Gothic}[AutoFakeSlant=0.15]\setmainfont{Malgun Gothic}[AutoFakeSlant=0.15]%
}{\IfFontExistsTF{Noto Sans CJK KR}{%
  \setsansfont{Noto Sans CJK KR}[AutoFakeSlant=0.15]\setmainfont{Noto Sans CJK KR}[AutoFakeSlant=0.15]%
}{%
  \setsansfont{NanumGothic}[AutoFakeSlant=0.15]\setmainfont{NanumGothic}[AutoFakeSlant=0.15]%
}}
\usetheme{sustech}

\usepackage{amssymb}
\usepackage{tikz}
\usetikzlibrary{positioning,arrows.meta,shapes}
\graphicspath{{figures/}}

% 조건 플래그(테마의 시각화 구분 페이지용, 사용 안 함)
\newif\ifvizbgdiv
\vizbgdivfalse

% --- 중국어 고정 문구를 한국어로 -------------------------------------------
\renewcommand{\figureprefix}{그림}
\renewenvironment{callout}[1][핵심]{\begin{alertblock}{#1}}{\end{alertblock}}
\makeatletter
\renewcommand{\titlemetaline}{%
  \ifdefempty{\presenter}{}{\presenter\metasep}%
  \ifdefempty{\presvenue}{}{\presvenue\metasep}%
  \insertdate}
\makeatother

% --- 브랜딩: 학교 로고·크레딧 끔 (채널 로고가 있으면 \setlogo[0.14\paperheight]{경로}) ---
\setlogo{}
\hidecredits
